\chapter*{Введение}
\addcontentsline{toc}{chapter}{Введение}
\label{ch:intro}

    Управление робототехническими системами по инструкциям на естественном языке является одной из центральных задач современной робототехники. За последние три года в этой области произошёл качественный сдвиг благодаря появлению визуально-языковых моделей действий (Vision-Language-Action, VLA) -- класса моделей, объединяющих визуальное восприятие, понимание языка и генерацию управляющих команд в единой архитектуре \cite{brohan2023rt2, kim2024openvla, black2024pi0}. Модели этого класса обучаются на крупномасштабных визуально-языковых данных из Интернета и дообучаются на робототехнических демонстрациях, что позволяет им следовать произвольным текстовым инструкциям и обобщать на новые объекты и среды.

    Интерес к VLA-моделям демонстрирует взрывной рост: по данным поиска на платформе OpenReview по ключевому слову <<Vision-Language-Action>>, на конференции ICLR 2024 была подана лишь одна такая работа, на ICLR 2025 -- девять, а к ICLR 2026 их число достигло 164, что соответствует примерно восемнадцатикратному росту всего за год \cite{reuss2025blog} (рисунок~\ref{fig:vla_growth}). Ключевыми представителями данного направления являются RT-2 \cite{brohan2023rt2} с 55 миллиардами параметров, показавший двукратное превосходство над предшественником на новых объектах; OpenVLA \cite{kim2024openvla} -- открытая модель с 7 миллиардами параметров, превосходящая RT-2-X на 16,5\% по метрике успешности; $\pi_0$ \cite{black2024pi0} от Physical Intelligence, использующая flow matching для генерации действий; и FLOWER \cite{reuss2025flower} -- компактная модель с 950 миллионами параметров, достигшая state-of-the-art на бенчмарке CALVIN.

\begin{figure}[htbp]
    \centering
    \begin{tikzpicture}
        \begin{axis}[
            thesisplot,
            ybar,
            bar width=26pt,
            height=5.6cm, width=0.66\textwidth,
            ylabel={Число работ по VLA},
            symbolic x coords={ICLR 2024, ICLR 2025, ICLR 2026},
            xtick=data,
            ymin=0, ymax=185,
            nodes near coords,
            nodes near coords style={font=\footnotesize},
            enlarge x limits=0.28,
        ]
        \addplot[fill=blue!45, draw=blue!60!black] coordinates {(ICLR 2024,1) (ICLR 2025,9) (ICLR 2026,164)};
        \end{axis}
    \end{tikzpicture}
    \caption{Рост числа работ по тематике Vision-Language-Action на ICLR (поиск по ключевому слову на OpenReview, по данным~\cite{reuss2025blog}).}
    \label{fig:vla_growth}
\end{figure}

    Однако, несмотря на впечатляющий прогресс, успешность выполнения операций VLA-моделями в реальных условиях остаётся на уровне 70--80\%, что недостаточно для надёжного практического применения \cite{reuss2025blog}. Между открытыми моделями и закрытыми фронтирными системами, такими как Gemini Robotics 1.5 \cite{deepmind2025gemini}, существует значительный разрыв в способности к zero-shot обобщению, который не виден из результатов на стандартных бенчмарках. Обучение с подкреплением (Reinforcement Learning, RL) рассматривается как одно из ключевых направлений преодоления этого разрыва. RL позволяет агенту улучшать политику управления через взаимодействие со средой на основе сигнала вознаграждения, что потенциально может довести успешность VLA-моделей с 70--80\% до требуемых 99\%+.

    Интеграция RL и визуально-языковых моделей развивается по нескольким направлениям. Первое -- RL-дообучение уже обученных VLA-моделей: работы PLD \cite{pld2026} (99\% на LIBERO) и stage-aware RL \cite{starevla2025} (98\% на SIMPLER), представленные на ICLR 2026, показывают, что residual RL и стадийные вознаграждения позволяют почти полностью закрыть разрыв на стандартных бенчмарках. Второе -- использование визуально-языковых моделей для автоматической генерации функций вознаграждения, заменяющих трудоёмкое ручное проектирование наград: EUREKA \cite{ma2024eureka} превосходит экспертные вознаграждения на 83\% задач, а Text2Reward \cite{xie2024text2reward} генерирует плотные reward-функции в виде исполняемого кода, не уступая экспертным на 13 из 17 задач. Третье -- применение VLM в качестве высокоуровневого планировщика в связке с RL-агентом \cite{ahn2022saycan, embodiedr1_2026}. Тем не менее, ни один из методов интеграции RL и VLA пока не утвердился как стандартный подход -- эта проблема остаётся открытой.

    \textbf{Объект исследования} -- процессы обучения и дообучения визуально-языковых моделей действий для робототехнического манипулирования с применением обучения с подкреплением.

    \textbf{Предмет исследования} -- методы RL-дообучения визуально-языковых моделей действий (residual RL, стадийные вознаграждения, заморозка базовой политики), закономерности прироста SR в зависимости от baseline-качества VLA, а также сравнительная эффективность VLM-генерации наград как альтернативного направления интеграции.

    \textbf{Актуальность.} Несмотря на высокие результаты VLA-моделей в симуляции (до 99\% на LIBERO), их надёжность в реальных условиях остаётся недостаточной. RL рассматривается как ключевой инструмент доведения SR до промышленно приемлемого уровня, однако отсутствует единая методика выбора схемы интеграции VLM/VLA и RL с учётом доступных данных и вычислительных ресурсов.

    \textbf{Научная новизна} работы состоит в следующем:
    \begin{enumerate}
        \item предложена единая классификация подходов интеграции RL и VLM/VLA по восьми критериям с количественным сопоставлением результатов из литературы;
        \item экспериментально сравнены девять открытых VLA-моделей (OpenVLA-OFT, UniVLA, $\pi_0$, CogACT, SmolVLA, $\pi_0$-fast, SpatialVLA, Octo, RT-1-X) на бенчмарке LIBERO в едином протоколе до и после RL-дообучения;
        \item установлено, что residual SAC со стадийным вознаграждением стабильно повышает SR на 1,5--4,1~п.п., тогда как end-to-end разморозка VLA приводит к деградации ($-$0,6~п.п.);
        \item показана закономерность: наибольший прирост от RL получают модели с более низким baseline SR (RT-1-X: +4,1~п.п.), наименьший -- у лидеров (OpenVLA: +1,5~п.п.);
        \item в дополнительной серии на MetaWorld подтверждено, что VLM-награды при 3 итерациях рефлексии поднимают SR на контактной задаче push-v3 с 48\% до 85\%.
    \end{enumerate}

    \textbf{Практическая значимость.} Основной результат -- выбор VLA и схемы RL-дообучения (residual SAC + стадийное вознаграждение при SR\,$>$\,90\%). Дополнительная серия VLM-наград покрывает сценарий без VLA (рисунок~\ref{fig:decision_tree}).

    Целью настоящей работы является разработка методических рекомендаций по применению обучения с подкреплением в задачах обучения и дообучения визуально-языковых моделей действий для робототехнического манипулирования. Для достижения цели поставлены следующие задачи:

\begin{enumerate}
    \item Провести аналитический обзор современных VLA-моделей и способов их интеграции с обучением с подкреплением на основе работ 2022--2026 гг.
    \item Систематизировать и классифицировать существующие подходы по трём направлениям: RL-дообучение VLA, VLM-генерация вознаграждений для RL, VLM как планировщик в связке с RL.
    \item Разработать методику проведения вычислительного эксперимента.
    \item Реализовать программный прототип RL-дообучения VLA и дополнительный модуль VLM-генерации наград в среде симуляции.
    \item Провести экспериментальное сравнение на стандартных бенчмарках с базовыми методами.
    \item Проанализировать результаты, выявить ограничения и сформулировать рекомендации по выбору архитектуры интеграции VLM и RL.
\end{enumerate}

    Работа организована следующим образом. В главе~\ref{ch:chap1} приведён аналитический обзор VLA-моделей, методов RL и подходов к их интеграции. В главе~\ref{ch:chap2} сформулирована постановка задачи, обоснован выбор RL-дообучения VLA, описана архитектура residual RL и методика эксперимента. В главе~\ref{ch:chap3} описана реализация прототипа: основные серии~1--4 (VLA+RL, LIBERO) и дополнительная серия~5 (VLM-награды, MetaWorld). В заключении сформулированы методические рекомендации.

    Центральная логика (рисунок~\ref{fig:thesis_roadmap}): \textbf{фокус} -- VLA-baseline $\to$ RL-схемы $\to$ RL всех моделей $\to$ абляция; \textbf{дополнение} -- VLM-награды на MetaWorld.

\begin{figure}[htbp]
    \centering
    \begin{tikzpicture}[
        node distance=0.85cm and 2.6cm,
        stage/.style={draw, rounded corners=2pt, align=center, text width=4.0cm, minimum height=1.0cm, font=\small},
        series/.style={draw, rounded corners=2pt, align=center, text width=4.4cm, minimum height=0.85cm, font=\small},
        line/.style={semithick},
        arr/.style={line, -{Stealth[length=3mm, width=2.2mm]}},
        dasharr/.style={line, dashed, -{Stealth[length=3mm, width=2.2mm]}}
    ]
        \node[stage, fill=green!12] (g1) {Глава 1:\\обзор VLA + RL};
        \node[stage, fill=blue!10, below=of g1] (g2) {Глава 2:\\постановка и методика};
        \node[stage, fill=orange!12, below=of g2] (g3) {Глава 3:\\VLA+RL (осн.) + VLM (доп.)};
        \node[stage, fill=purple!10, below=of g3] (g4) {Заключение:\\рекомендации};

        \draw[arr] (g1.south) -- (g2.north);
        \draw[arr] (g2.south) -- (g3.north);
        \draw[arr] (g3.south) -- (g4.north);

        \node[series, fill=gray!10, right=2.8cm of g2] (s1) {С1: VLA-baseline (9 моделей)};
        \node[series, fill=gray!10, below=0.45cm of s1] (s2) {С2: RL-схемы (OpenVLA)};
        \node[series, fill=gray!10, below=0.45cm of s2] (s3) {С3: RL всех VLA (residual SAC)};
        \node[series, fill=gray!10, below=0.45cm of s3] (s4) {С4: абляция (stage, freeze)};
        \node[series, fill=orange!10, below=0.45cm of s4] (s5) {С5: VLM-награда (MetaWorld)};

        \node[draw, dashed, rounded corners=3pt, inner sep=0.35cm, fit=(s1)(s5), label={[font=\small]above:Экспериментальные серии}] {};

        \coordinate (linkH) at ($(g3.east)+(0.9,0)$);
        \coordinate (linkV) at (linkH |- s2.west);
        \draw[line, dashed] (g3.east) -- (linkH);
        \draw[line, dashed] (linkH) -- (linkV);
        \draw[dasharr] (linkV) -- (s2.west);
    \end{tikzpicture}
    \caption{Структура работы и логика экспериментальных серий.}
    \label{fig:thesis_roadmap}
\end{figure}

\endinput
